Une fonction est d'abord une machine à convertir : au collège, un coefficient de
proportionnalité, une échelle, un pourcentage. Elle devient un objet d'étude au lycée, puis
le support d'une théorie — l'analyse — qui occupe la moitié du programme de terminale.

La progression est celle du regard. On regarde d'abord une fonction en un point : son image,
son antécédent. Puis sur un intervalle : son sens de variation, ses extremums, sa forme
canonique s'il s'agit d'un trinôme. Puis à l'infini, avec les suites, où apparaît pour la
première fois une notion qui ne se calcule pas mais s'approche : la limite. Le reste du thème
consiste à équiper cette notion — limites de fonctions, continuité, dérivation — puis à
l'employer : exponentielle, logarithme, trigonométrie, et enfin l'intégrale.

C'est le thème où le programme \emph{admet} le plus, et donc celui où la formalisation
déplace le plus de choses. Trois résultats en particulier — la convergence des suites
croissantes majorées, le théorème des valeurs intermédiaires, l'égalité entre l'intégrale et
l'aire sous la courbe — reposent tous, en dernière analyse, sur la complétude de $\mathbb{R}$,
qui n'est jamais énoncée au lycée.
